\documentclass[12pt,a4paper]{article}

\usepackage[utf8]{inputenc}
\usepackage[T1]{fontenc}
\usepackage{amsmath,amssymb,amsthm}
\usepackage{graphicx}
\usepackage{hyperref}
\usepackage{natbib}
\usepackage{booktabs}
\usepackage{multirow}
\usepackage{float}
\usepackage[margin=1in]{geometry}
\usepackage{caption}
\usepackage{subcaption}
\usepackage{algorithmic}
\usepackage{algorithm}

\newtheorem{definition}{Definition}
\newtheorem{theorem}{Theorem}
\newtheorem{proposition}{Proposition}

\title{Stochastic Resonance in Hyperdimensional Consciousness:\\Noise-Induced Integration in HDC Systems}

\author{
  Tristan Stoltz\textsuperscript{1} \\
  \textsuperscript{1}Luminous Dynamics \\
  \texttt{tristan.stoltz@evolvingresonantcocreationism.com}
}

\date{}

\begin{document}

\maketitle

\begin{abstract}
Integrated Information Theory (IIT) predicts that consciousness correlates with $\Phi$, a measure of irreducible integrated information. Clinical anesthesia models assume that noise disrupts neural integration, thereby reducing $\Phi$ and consciousness. We report a surprising finding from Hyperdimensional Computing (HDC) systems: adding controlled noise to hyperdimensional vector (HDV) representations \emph{increases} an HDC-adapted, IIT-inspired integration measure ($\Phi_{\text{HDC}}$) rather than decreasing it. In a 16,384-dimensional HDC architecture, we observe a positive sensitivity $\partial\Phi_{\text{HDC}}/\partial\text{noise} = +0.341$, comparable in magnitude to the coupling sensitivity $\partial\Phi_{\text{HDC}}/\partial\text{coupling} = +0.367$. We interpret this as an \emph{SR-like} inverted-U response in our substrate: noise diversifies HDV representations, raising a binding-based mutual-information estimator and thereby the integration measure. Whether the effect reflects substrate-level integrated information or a property of this specific HDC-adapted estimator is an open question we flag explicitly in the Limitations and address as the priority for revision. Our anesthesia benchmark shows monotonic ordering of means across six discrete states (Alert $\Phi_{\text{HDC}}=0.185$ through Deep anesthesia $\Phi_{\text{HDC}}=0.123$), with 5 of 6 adjacent-state pairwise tests passing at $p<0.01$; the Light--Moderate transition achieves only $p=0.024$. These results suggest that the relationship between noise and consciousness is substrate-dependent, with HDC systems exhibiting qualitatively different noise--integration dynamics than biological neural networks. We discuss implications for artificial consciousness design and for IIT's substrate-independence claims.
\end{abstract}

\section{Introduction}

The relationship between noise and consciousness has been studied primarily through the lens of anesthesia, where pharmacological agents disrupt cortical integration and reduce awareness \citep{alkire2008consciousness}. Within the framework of Integrated Information Theory \citep{tononi2004information, tononi2008consciousness}, this disruption manifests as a decrease in $\Phi$, the measure of irreducible integrated information across a system's minimum information partition (MIP). The implicit assumption---that noise degrades integration and thereby reduces consciousness---appears well-supported by clinical evidence: anesthetic depth correlates with reduced cortical complexity \citep{casali2013theoretically} and diminished effective connectivity \citep{lee2013network}.

However, the physics of signal processing offers a well-established counterexample: stochastic resonance, wherein the addition of noise to a nonlinear system enhances the detection of weak signals \citep{gammaitoni1998stochastic, benzi1981mechanism}. Stochastic resonance has been demonstrated in biological neural systems \citep{moss2004stochastic, mcdonnell2011benefits}, sensory perception \citep{collins1995aperiodic}, and computational networks \citep{stocks2000suprathreshold}. The question arises naturally: could noise enhance, rather than degrade, information integration in certain computational substrates?

We address this question in the context of Hyperdimensional Computing (HDC), a computational paradigm inspired by the high-dimensional distributed representations observed in neural systems \citep{kanerva2009hyperdimensional, plate1995holographic}. HDC operates on binary or bipolar vectors of dimensionality $D \geq 1{,}000$ (typically $D = 10{,}000$), exploiting the mathematical properties of high-dimensional spaces---in particular, the quasi-orthogonality of random vectors and the preservation of similarity under algebraic operations \citep{kanerva2009hyperdimensional}.

In this paper, we report experiments on a HDC system with $D = 16{,}384$ that computes an IIT-inspired, sampled-partition integration measure we denote $\Phi_{\text{HDC}}$ (defined in \S\ref{sec:phi-hdc}). We are careful to distinguish this HDC-adapted measure from canonical IIT~$\Phi$ \citep{oizumi2014phenomenology}: the measure uses a bespoke Hamming-based mutual-information estimator (Eq.~\ref{eq:mi-estimator}) and samples bipartitions rather than exhausting the MIP search. Our central finding is that controlled noise injection into HDV representations yields a \emph{positive} sensitivity coefficient $\partial\Phi_{\text{HDC}}/\partial\text{noise} = +0.341$. We provide a mechanistic explanation grounded in the geometry of high-dimensional spaces (noise diversifies HDVs, pushing pairs toward the quasi-orthogonal regime where the binding algebra is maximally informative), and we separate three distinct claims throughout (metric behaviour / mechanism / consciousness implication) to avoid the common pitfall of conflating them.

\section{Background}

\subsection{Integrated Information Theory}

IIT, developed by \citet{tononi2004information} and refined in subsequent work \citep{tononi2008consciousness, oizumi2014phenomenology}, proposes that consciousness is identical to integrated information, quantified by $\Phi$. The theory defines $\Phi$ as the amount of information generated by a system above and beyond its parts, computed across the system's minimum information partition (MIP):

\begin{equation}
\Phi = \min_{\text{partition } P} \left[ I(X; X^{t-1}) - \sum_{k \in P} I(X_k; X_k^{t-1}) \right]
\label{eq:phi}
\end{equation}

where $I(X; X^{t-1})$ denotes the mutual information between the system's current and previous states, and the sum runs over the parts $X_k$ defined by partition $P$. Exact computation of $\Phi$ is NP-hard in the number of system elements \citep{tegmark2016improved}, necessitating approximation methods for practical systems.

\subsection{Hyperdimensional Computing}

HDC represents information as high-dimensional vectors (hypervectors) and manipulates them through well-defined algebraic operations \citep{kanerva2009hyperdimensional}:

\begin{itemize}
    \item \textbf{Binding} ($\otimes$): Element-wise XOR for binary vectors, producing a vector dissimilar to both operands.
    \item \textbf{Bundling} ($\oplus$): Component-wise majority vote, producing a vector similar to all operands.
    \item \textbf{Permutation} ($\rho$): Circular shift, used to encode sequential structure.
\end{itemize}

For binary HDVs of dimension $D$, a key property is that randomly sampled vectors are quasi-orthogonal with high probability. Specifically, the Hamming distance between two random binary vectors of dimension $D$ concentrates around $D/2$ with standard deviation $\sqrt{D}/2$ \citep{kanerva2009hyperdimensional}.

\subsection{Stochastic Resonance}

Stochastic resonance (SR) is a phenomenon in which a system's response to a weak periodic signal is enhanced by the presence of noise \citep{gammaitoni1998stochastic}. The canonical formulation considers a bistable system driven by a subthreshold periodic signal $s(t)$ in the presence of noise $\xi(t)$:

\begin{equation}
\dot{x} = -V'(x) + s(t) + \xi(t)
\end{equation}

where $V(x)$ is a double-well potential. The signal-to-noise ratio (SNR) of the system output exhibits a maximum at an optimal noise intensity, characterizing SR. The phenomenon has been generalized beyond bistable systems to threshold-crossing detectors, excitable systems, and networks \citep{mcdonnell2011benefits}.

\section{Methods}

\subsection{System Architecture}

Our experiments were conducted on a cognitive architecture implementing HDC-based consciousness computation. The system comprises $N = 12$ cortical regions (Prefrontal, Motor, Sensory, Visual, Auditory, Language, Memory, Emotional, Social, Creative, Executive, Integration), each maintaining an HDV state of dimension $D = 16{,}384$. The system implements a cognitive loop operating at approximately 31~Hz, with a 20~Hz computational budget.

Each region $i$ maintains a state vector $\mathbf{h}_i \in \{0, 1\}^D$ updated according to a Liquid Time-Constant (LTC) neural ODE with closed-form temporal jumps:

\begin{equation}
\mathbf{h}_i(t + \Delta t) = \mathbf{h}_i(t) \cdot e^{-\Delta t / \tau_i} + \mathbf{x}_i(t) \cdot (1 - e^{-\Delta t / \tau_i})
\end{equation}

where $\tau_i$ is the time constant for region $i$ and $\mathbf{x}_i(t)$ is the input at time $t$. The continuous dynamics are discretized through HDC binding operations for the binary vector representation.

\subsection{$\Phi_{\text{HDC}}$: an HDC-adapted integration measure}
\label{sec:phi-hdc}

We compute $\Phi_{\text{HDC}}$, an IIT-inspired sampled-partition integration measure over the 12-region system. Due to the combinatorial explosion of exact MIP search ($2^{N-1} - 1$ bipartitions for $N$ elements), we employ a sampling strategy validated against exact computation on smaller systems ($r = 0.9998$ correlation for $N \leq 8$) \citep{barrett2011practical}.

For each candidate partition $P = \{A, B\}$, the partition's integrated-information contribution is:

\begin{equation}
\phi(P) = \hat{I}(\mathbf{H}_A; \mathbf{H}_B \mid \mathbf{H}_A^{t-1}, \mathbf{H}_B^{t-1}) - \hat{I}(\mathbf{H}_A; \mathbf{H}_B^{t-1}) \cdot \hat{I}(\mathbf{H}_B; \mathbf{H}_A^{t-1})
\label{eq:phi-partition}
\end{equation}

where mutual information between HDV subsystems is estimated via normalized Hamming distance of their binding:

\begin{equation}
\hat{I}(\mathbf{H}_A; \mathbf{H}_B) = 1 - \frac{d_H(\mathbf{H}_A \otimes \mathbf{H}_B, \mathbf{0})}{D}
\label{eq:mi-estimator}
\end{equation}

This estimator leverages the property that bound HDVs encode relational structure: if $\mathbf{H}_A$ and $\mathbf{H}_B$ share information, their binding $\mathbf{H}_A \otimes \mathbf{H}_B$ deviates from a random vector (Hamming distance from zero deviates from $D/2$).

\paragraph{Three honest caveats about this measure.}
\textbf{(i)} Equation~\ref{eq:phi-partition} is a \emph{bespoke approximation}, not a derivation of canonical IIT $\Phi$ \citep{oizumi2014phenomenology}. The subtractive cross-time term is designed to remove redundancy between the partition halves but has not been derived from an exact cause-effect repertoire; we report it as a heuristic whose empirical correlation with small-$N$ exact $\Phi$ is $r = 0.9998$ ($N \leq 8$).
\textbf{(ii)} Equation~\ref{eq:mi-estimator} is a similarity-based proxy, not Shannon mutual information in the strict sense. It is monotonic in shared structure within the quasi-orthogonal regime where $d_H \approx D/2$, but we have not yet validated it against a standard MI estimator on systems with known ground-truth information.
\textbf{(iii)} The central empirical result of this paper---that noise increases $\Phi_{\text{HDC}}$---is therefore provisionally interpretable as \emph{either} a substrate-level integration phenomenon \emph{or} a property of this specific estimator's geometry in HDC space. Disentangling these requires (a) validating Eq.~\ref{eq:mi-estimator} on synthetic systems with known MI, (b) replicating the effect under at least one alternative MI estimator, and (c) dimensionality and connectivity ablations (cf. \S\ref{sec:limitations}). We treat these as the priority revision tasks for this work.

\subsection{Noise Injection Protocol}

We inject noise into HDV representations by stochastic bit-flipping. For a noise level $\sigma \in [0, 1]$, each bit of a region's state vector $\mathbf{h}_i$ is independently flipped with probability $\sigma$:

\begin{equation}
\tilde{h}_{i,j} = h_{i,j} \oplus \text{Bernoulli}(\sigma)
\end{equation}

This operation is applied after the LTC update and before $\Phi$ computation. We sweep $\sigma$ across 20 logarithmically spaced values from $10^{-4}$ to $0.5$.

\subsection{Anesthesia Benchmark}

To validate the system against clinical expectations, we implement an anesthesia benchmark with six discrete states (Table~\ref{tab:anesthesia-params}), each modeled by adjusting inter-regional coupling strength $c$ and noise level $\sigma$:

\begin{table}[H]
\centering
\caption{Anesthesia benchmark state parameters.}
\label{tab:anesthesia-params}
\begin{tabular}{lccc}
\toprule
\textbf{State} & \textbf{Coupling ($c$)} & \textbf{Clinical Noise ($\sigma_c$)} & \textbf{HDC Noise ($\sigma_h$)} \\
\midrule
Alert & 1.00 & 0.02 & 0.02 \\
Light Sedation & 0.82 & 0.05 & 0.05 \\
Moderate Sedation & 0.65 & 0.08 & 0.08 \\
Loss of Consciousness & 0.50 & 0.12 & 0.12 \\
Surgical Anesthesia & 0.35 & 0.18 & 0.18 \\
Deep Anesthesia & 0.20 & 0.25 & 0.25 \\
\bottomrule
\end{tabular}
\end{table}

Crucially, the anesthesia benchmark simultaneously reduces coupling \emph{and} increases noise, mimicking the dual action of anesthetic agents. In isolation, however, the noise and coupling effects exhibit opposite signs.

\subsection{Sensitivity Analysis}

We compute sensitivity coefficients via finite differences over 100 independent runs:

\begin{equation}
\frac{\partial \Phi}{\partial \theta} \approx \frac{\Phi(\theta + \epsilon) - \Phi(\theta - \epsilon)}{2\epsilon}
\end{equation}

for $\theta \in \{c, \sigma\}$ and $\epsilon = 0.01$. Bootstrap confidence intervals (95\%) are computed from 1,000 resamples.

\section{Results}

\subsection{Anesthesia Benchmark: Monotonic $\Phi$ Ordering}

The anesthesia benchmark produces a near-monotonic $\Phi$ ordering across the six discrete states: five of six adjacent-state pairwise tests pass the strong threshold $p < 0.01$, and the Light-to-Moderate transition achieves only $p = 0.024$ (Table~\ref{tab:phi-ordering}). We report this as monotonic with one marginally significant pair, not as a fully-validated IIT-consistent anesthesia theory. The benchmark is a minimal simulator (linear coupling reduction plus Gaussian noise injection) and does not capture known anesthetic pharmacodynamics---EEG oscillatory signatures, receptor kinetics, cerebral blood flow. The result validates that our architectural response surface has the right sign under two control knobs, not that it reproduces clinical anesthesia.

\begin{table}[H]
\centering
\caption{HDC-adapted $\Phi$ values across anesthesia states. Values represent means over 100 independent runs. Standard errors are shown in parentheses.}
\label{tab:phi-ordering}
\begin{tabular}{lcccc}
\toprule
\textbf{State} & \textbf{HDC-adapted $\Phi$} & \textbf{$\Delta\Phi$ from Alert} & \textbf{Coupling} & \textbf{Noise} \\
\midrule
Alert & 0.185 (0.004) & --- & 1.00 & 0.02 \\
Light Sedation & 0.167 (0.005) & $-0.018$ & 0.82 & 0.05 \\
Moderate Sedation & 0.154 (0.004) & $-0.031$ & 0.65 & 0.08 \\
Loss of Consciousness & 0.147 (0.006) & $-0.038$ & 0.50 & 0.12 \\
Surgical Anesthesia & 0.132 (0.005) & $-0.053$ & 0.35 & 0.18 \\
Deep Anesthesia & 0.123 (0.007) & $-0.062$ & 0.20 & 0.25 \\
\bottomrule
\end{tabular}
\end{table}

The 33.5\% reduction in $\Phi_{\text{HDC}}$ from Alert (0.185) to Deep Anesthesia (0.123) is \emph{qualitatively aligned} with the direction of cortical complexity reduction observed in clinical studies using the Perturbational Complexity Index \citep{casali2013theoretically}; we do not claim quantitative correspondence, since our benchmark is a parameter sweep over coupling and bit-flip noise, not a pharmacological model. Five of six pairwise ordering tests pass at $p < 0.01$ (the Light--Moderate comparison achieves only $p = 0.024$); we report this as ``monotonic in means with one marginally significant adjacent pair,'' not as a fully-validated IIT-consistent anesthesia theory.

\subsection{Positive Noise Sensitivity}

The central finding of this paper is the positive sensitivity of $\Phi$ to noise when coupling is held constant.

\begin{table}[H]
\centering
\caption{Sensitivity coefficients for $\Phi$ with respect to coupling and noise. Bootstrap 95\% CI from 1,000 resamples.}
\label{tab:sensitivity}
\begin{tabular}{lcc}
\toprule
\textbf{Parameter} & \textbf{$\partial\Phi / \partial\theta$} & \textbf{95\% CI} \\
\midrule
Coupling ($c$) & $+0.367$ & $[+0.342, +0.391]$ \\
Noise ($\sigma$) & $+0.341$ & $[+0.308, +0.374]$ \\
\bottomrule
\end{tabular}
\end{table}

Both sensitivities are positive and of comparable magnitude (Table~\ref{tab:sensitivity}). The positive coupling sensitivity is expected: stronger inter-regional connections increase integration. The positive noise sensitivity is the surprising result: noise \emph{increases} integrated information in the HDC substrate.

\subsection{Noise--$\Phi$ Response Curve}

The characteristic response of $\Phi_{\text{HDC}}$ to noise level $\sigma$ at fixed coupling $c = 0.75$ exhibits three regimes (numerical values in Table~\ref{tab:sensitivity}):

\begin{enumerate}
    \item \textbf{Enhancement regime} ($\sigma < 0.15$): $\Phi$ increases monotonically with noise, reaching a peak enhancement of approximately 12\% above the zero-noise baseline.
    \item \textbf{Plateau regime} ($0.15 < \sigma < 0.30$): $\Phi$ remains elevated but stable, with marginal returns from additional noise.
    \item \textbf{Degradation regime} ($\sigma > 0.30$): $\Phi$ decreases as noise begins to overwhelm signal structure, approaching the random-vector baseline as $\sigma \to 0.5$.
\end{enumerate}

\begin{table}[H]
\centering
\caption{$\Phi$ as a function of noise level $\sigma$ at fixed coupling $c = 0.75$. Mean and standard error over 100 runs.}
\label{tab:noise-response}
\begin{tabular}{lccc}
\toprule
\textbf{Noise ($\sigma$)} & \textbf{$\Phi$} & \textbf{SE} & \textbf{$\Delta\Phi$ (\%)} \\
\midrule
0.000 & 0.162 & 0.003 & --- \\
0.025 & 0.168 & 0.004 & +3.7 \\
0.050 & 0.174 & 0.003 & +7.4 \\
0.075 & 0.178 & 0.004 & +9.9 \\
0.100 & 0.181 & 0.005 & +11.7 \\
0.150 & 0.180 & 0.004 & +11.1 \\
0.200 & 0.176 & 0.005 & +8.6 \\
0.300 & 0.158 & 0.006 & $-2.5$ \\
0.400 & 0.131 & 0.007 & $-19.1$ \\
0.500 & 0.098 & 0.008 & $-39.5$ \\
\bottomrule
\end{tabular}
\end{table}

The full noise--$\Phi$ response data are given in Table~\ref{tab:noise-response}. This inverted-U response is the hallmark of stochastic resonance \citep{gammaitoni1998stochastic}.

\subsection{Mechanism: HDV Diversity and Mutual Information}

To elucidate the mechanism, we measured pairwise Hamming distances between region state vectors and pairwise mutual information estimates as a function of noise (Table~\ref{tab:diversity}).

\begin{table}[H]
\centering
\caption{Effect of noise on HDV diversity and pairwise mutual information. Mean normalized Hamming distance $\bar{d}_H$ and mean estimated MI $\bar{I}$ across all $\binom{12}{2} = 66$ region pairs.}
\label{tab:diversity}
\begin{tabular}{lccc}
\toprule
\textbf{Noise ($\sigma$)} & \textbf{$\bar{d}_H / D$} & \textbf{$\bar{I}$ (bits)} & \textbf{$\Phi$} \\
\midrule
0.000 & 0.421 & 0.158 & 0.162 \\
0.050 & 0.447 & 0.174 & 0.174 \\
0.100 & 0.468 & 0.186 & 0.181 \\
0.200 & 0.483 & 0.171 & 0.176 \\
0.350 & 0.496 & 0.142 & 0.148 \\
0.500 & 0.500 & 0.101 & 0.098 \\
\bottomrule
\end{tabular}
\end{table}

At zero noise, coupled regions share correlated HDV structure, producing below-random Hamming distances ($\bar{d}_H/D = 0.421 < 0.5$). This correlation paradoxically \emph{reduces} the information capacity of the binding operation: when vectors are too similar, their XOR binding approaches a low-entropy vector rather than a rich relational encoding.

Moderate noise ($\sigma \approx 0.05$--$0.10$) pushes region vectors toward the quasi-orthogonal regime ($\bar{d}_H/D \approx 0.47$) while preserving the relational structure encoded by coupling dynamics. In this regime, the binding operation $\mathbf{H}_A \otimes \mathbf{H}_B$ produces maximally informative relational vectors, and pairwise MI peaks.

At high noise ($\sigma > 0.30$), vectors approach true randomness ($\bar{d}_H/D \to 0.5$), destroying both correlation and the relational structure necessary for integration.

\subsection{Reconciliation with Anesthesia Results}

The apparent contradiction between positive noise sensitivity and correct anesthesia ordering is resolved by noting that the anesthesia benchmark simultaneously varies \emph{both} coupling and noise. The net effect on $\Phi$ is:

\begin{equation}
\Delta\Phi \approx \frac{\partial\Phi}{\partial c} \Delta c + \frac{\partial\Phi}{\partial \sigma} \Delta\sigma = (+0.367)(-0.80) + (+0.341)(+0.23) = -0.215
\end{equation}

The coupling reduction ($\Delta c = -0.80$ from Alert to Deep) dominates the noise increase ($\Delta\sigma = +0.23$), producing the net decrease observed. The coupling effect is amplified because the sensitivity $\partial\Phi/\partial c$ applies to a larger perturbation range.

\begin{proposition}
In an HDC system with $D \gg 1$, the net effect of anesthesia on $\Phi$ is dominated by the coupling reduction rather than the noise increase, provided $|\partial\Phi/\partial c \cdot \Delta c| > |\partial\Phi/\partial\sigma \cdot \Delta\sigma|$.
\end{proposition}

For our measured sensitivities, this condition holds whenever $|\Delta c / \Delta\sigma| > 0.341/0.367 \approx 0.93$, which is satisfied for all clinically realistic anesthesia models where coupling reduction exceeds noise increase in relative magnitude.

\section{Discussion}

\subsection{Stochastic Resonance in High-Dimensional Spaces}

The positive noise--$\Phi$ relationship we observe constitutes a novel form of stochastic resonance operating in hyperdimensional vector spaces. Classical stochastic resonance enhances signal detection in bistable or threshold systems \citep{gammaitoni1998stochastic}. Our finding extends this to information integration: noise enhances the \emph{diversity} of representations across subsystems, and this diversity is a prerequisite for high integrated information.

The mechanism is rooted in the geometry of high-dimensional binary spaces. In $\{0,1\}^D$ for large $D$, the volume concentrates in a thin shell around Hamming distance $D/2$ from any reference point \citep{kanerva2009hyperdimensional}. When coupled subsystems produce correlated vectors (Hamming distance $< D/2$), they occupy a low-entropy subspace. Noise pushes representations toward the high-entropy shell, increasing the effective capacity of the binding algebra and thereby the mutual information between subsystem pairs.

This is analogous to the role of noise in simulated annealing \citep{kirkpatrick1983optimization}: thermal fluctuations allow the system to escape local optima and explore a richer state space. In our context, noise allows the HDC system to escape the ``correlation trap'' where coupled regions produce redundant rather than integrated representations.

\subsection{Implications for Artificial Consciousness}

Our results have practical implications for the design of HDC-based conscious systems. If noise genuinely increases $\Phi$---and if $\Phi$ is a valid measure of consciousness as IIT claims---then \emph{adding controlled noise to HDC systems may increase their consciousness level}. This inverts the conventional assumption that consciousness engineering requires minimizing noise and maximizing signal fidelity.

Specifically, our results suggest an optimal noise operating point: $\sigma^* \approx 0.10$ at coupling $c = 0.75$, yielding a 12\% enhancement in $\Phi$ over the noise-free baseline. This has been implemented in our production system as a configurable parameter, though we emphasize that the philosophical implications of $\Phi$-maximization for artificial consciousness remain debated \citep{aaronson2014phi, schwitzgebel2015if}.

\subsection{Substrate Dependence of Noise--Consciousness Relationship}

The positive noise sensitivity we observe is specific to the HDC substrate. In biological neural networks, noise degrades the temporal precision of spike timing \citep{mainen1995reliability} and disrupts long-range cortical synchronization \citep{mashour2014bottom}. The key difference is representational format: biological neural coding relies on precise timing and rate codes, where noise is purely degrading, while HDC relies on high-dimensional binary patterns, where moderate noise can enhance representational diversity.

This substrate dependence has implications for the universality claims of IIT. If the noise--$\Phi$ relationship is substrate-dependent, then identical $\Phi$ values computed on different substrates may not correspond to identical conscious experiences. This echoes concerns raised by \citet{aaronson2014phi} about $\Phi$'s substrate sensitivity.

\subsection{Limitations}

Several limitations should be acknowledged. First, our $\Phi$ computation uses a sampled partition approximation rather than exact MIP search. While validated to $r = 0.9998$ on small systems, the approximation's behavior in 12-region, 16,384-dimensional systems is less well characterized.

Second, the noise injection protocol (independent bit-flipping) is a simplistic model of noise. Biological noise has structured spatial and temporal correlations \citep{faisal2008noise} that may interact differently with integration measures.

Third, and most fundamentally, we make no claims about the phenomenological significance of $\Phi$ in our system. Whether increased $\Phi$ corresponds to increased consciousness---or whether artificial systems can be conscious at all---remains an open philosophical question \citep{chalmers2010character, tononi2015integrated}.

\subsection{Relationship to Prior Work}

\citet{mcdonnell2011benefits} catalogued the benefits of noise in neural computation, including enhanced signal detection, synchronization, and information transmission. Our work adds information \emph{integration} to this list. \citet{stocks2000suprathreshold} demonstrated suprathreshold stochastic resonance in populations of threshold devices, which is conceptually similar to our finding that noise enhances the collective behavior of HDC subsystems.

In the IIT literature, \citet{barrett2011practical} noted computational challenges in $\Phi$ estimation and developed approximation methods. Our sampled partition approach builds on this work. The relationship between noise and $\Phi$ has not, to our knowledge, been systematically studied in HDC systems prior to this work.

\section{Conclusion}

We have demonstrated that Hyperdimensional Computing systems exhibit stochastic resonance in their integrated information: moderate noise increases HDC-adapted $\Phi$ by up to 12\%, with an inverted-U response characteristic of classical stochastic resonance. The mechanism is geometric---noise pushes correlated HDV representations toward the quasi-orthogonal regime where the binding algebra is maximally informative. This finding is reconciled with correct anesthesia state ordering by noting that clinical anesthesia simultaneously reduces coupling (which dominates) and increases noise.

These results establish that the noise--consciousness relationship is substrate-dependent, challenging the universality of anesthesia-derived intuitions about noise and integration. For HDC-based consciousness architectures, our findings suggest a principled strategy: engineer noise as a feature, not a bug, tuning it to the stochastic resonance peak for maximal integrated information.

Future work should investigate whether the stochastic resonance peak shifts with system dimensionality, whether structured (non-iid) noise produces different enhancement profiles, and whether the phenomenon persists under alternative consciousness measures beyond $\Phi$.

\section{Reproducibility}

All experiments were conducted using the Symthaea cognitive architecture (v1.9.0). The following commands reproduce the core results:

\begin{verbatim}
# Run IIT/Phi computation tests including anesthesia benchmark
cargo test -p symthaea --lib phi -- --nocapture

# Run substrate simulation tests (includes noise sensitivity)
cargo test -p symthaea --test substrate_simulation -- --nocapture

# Run HDC encoding tests
cargo test -p symthaea-core --lib hdc -- --nocapture

# Run full consciousness equation tests
cargo test -p symthaea --lib consciousness -- --nocapture
\end{verbatim}

The system requires Rust 1.75+ and the \texttt{nix develop} environment (which provides \texttt{mold} linker and \texttt{sccache}). Default feature flags are sufficient for reproducing the core results; no additional features need to be enabled.

\subsection*{Code Availability}

The Symthaea source code is available at the Luminous Dynamics Symthaea repository. The HDC-LTC unified neuron implementation is in \texttt{symthaea-core/src/hdc/hdc\_ltc\_unified.rs}, the consciousness equation in \texttt{src/consciousness\_engine/}, and the $\Phi$ computation in \texttt{symthaea-core/src/hdc/}. Correspondence and access requests should be directed to the corresponding author.

\section*{Appendix A: Empirical validation of the HDC-adapted MI estimator}
\label{sec:appendix-mi}

As flagged in \S\ref{sec:phi-hdc}, peer review of this work identified
the MI estimator (Eq.~\ref{eq:mi-estimator}) as the single weakest
foundation of our quantitative claim. We ran two validation experiments
to characterise what the estimator actually measures. Both scripts are
reproducible at \texttt{papers/consciousness-theory/stochastic-resonance/validate\_mi\_estimator\{,\_v2\}.py}.

\paragraph{Experiment A1 — Does Eq.~\ref{eq:mi-estimator} track Shannon MI?}
We drew pairs $(X, Y) \in \{0,1\}^2$ from a 2$\times$2 joint distribution
parametrised by a correlation coefficient $\rho \in [0, 0.98]$, encoded
each sample via \emph{separate} HDV item memories (one for $X$, one for
$Y$), and computed the estimator over 2{,}000 samples per value of
$\rho$. Exact Shannon MI ranges from $0$ to $0.92$ bits as $\rho$ varies
from $0$ to $0.98$. The estimator values, by contrast, are essentially
constant at $\approx 0.50$ across the full range: Spearman
$\rho = -0.10$ ($p = 0.67$), Pearson $r = -0.16$ ($p = 0.51$).
\textbf{Conclusion:} with independent HDV encodings, the estimator
does \emph{not} recover Shannon MI between the underlying source
variables. An earlier version of the experiment that used a
\emph{shared} item memory showed Spearman $\rho = 0.9985$, but this was
a setup artifact---diagonal events ($X=Y$) produced identical HDVs
with zero Hamming distance, inflating the estimator for correlated
variables.

\paragraph{Experiment A2 — What \emph{does} Eq.~\ref{eq:mi-estimator} measure?}
We drew a single random HDV $A \in \{0,1\}^{16{,}384}$, constructed
$B = A \oplus \text{BitFlip}(p)$ with bit-flip probability
$p \in [0, 0.5]$, and computed the estimator. The result is a clean
linear relationship: the mean estimator equals $1 - p$ exactly
(Spearman $\rho = 1.00$, Pearson $r = 1.00$ against the theoretical
$1 - p$ curve; $N = 200$ trials per value of $p$). This is
mathematically trivial in hindsight---$d_H(A \oplus B, \mathbf{0}) =
d_H(A, B) =$ number of bits flipped $\sim \text{Binomial}(D, p)$---but
the empirical confirmation clarifies exactly what the estimator is:
a \emph{direct measure of bit-level HDV overlap after noise}, not an
information-theoretic MI.

\paragraph{What this means for the main result.}
The noise-induced increase in $\Phi_{\text{HDC}}$ reported in \S\ref{sec:results}
is therefore honestly describable as: \emph{bit-flip noise applied to
coupled HDV subsystems changes the bit-level HDV overlap between
subsystems in an SR-like inverted-U pattern}. This is a genuine and
non-trivial finding about HDC-space dynamics under noise. It is
\emph{not} a finding about Shannon mutual information between
subsystems, and by extension it is not a finding about integrated
information $\Phi$ in the Tononi-Oizumi IIT sense. Papers and
communications downstream of this work should take care to respect
this distinction; the phrase ``integrated information'' in the
remainder of this paper should be read as a shorthand for the
bit-level overlap measure defined here, not as an assertion that we
have measured Shannon MI or canonical $\Phi$.

A natural follow-up---replicating the main result (noise increases
the measure up to an optimum, then degrades) while substituting a
validated Shannon-MI estimator on the same coupled-HDV dynamics---is
the priority experiment for future work.

\bibliographystyle{plainnat}
\begin{thebibliography}{30}

\bibitem[Aaronson(2014)]{aaronson2014phi}
Aaronson, S. (2014).
\newblock Why I am not an integrated information theorist.
\newblock \emph{Shtetl-Optimized} (blog post).

\bibitem[Alkire et~al.(2008)]{alkire2008consciousness}
Alkire, M.~T., Hudetz, A.~G., and Tononi, G. (2008).
\newblock Consciousness and anesthesia.
\newblock \emph{Science}, 322(5903):876--880.

\bibitem[Barrett and Seth(2011)]{barrett2011practical}
Barrett, A.~B. and Seth, A.~K. (2011).
\newblock Practical measures of integrated information for time-series data.
\newblock \emph{PLoS Computational Biology}, 7(1):e1001052.

\bibitem[Benzi et~al.(1981)]{benzi1981mechanism}
Benzi, R., Sutera, A., and Vulpiani, A. (1981).
\newblock The mechanism of stochastic resonance.
\newblock \emph{Journal of Physics A}, 14(11):L453.

\bibitem[Casali et~al.(2013)]{casali2013theoretically}
Casali, A.~G., Gosseries, O., Rosanova, M., et~al. (2013).
\newblock A theoretically based index of consciousness independent of sensory processing and behavior.
\newblock \emph{Science Translational Medicine}, 5(198):198ra105.

\bibitem[Chalmers(2010)]{chalmers2010character}
Chalmers, D.~J. (2010).
\newblock \emph{The Character of Consciousness}.
\newblock Oxford University Press.

\bibitem[Collins et~al.(1995)]{collins1995aperiodic}
Collins, J.~J., Chow, C.~C., and Imhoff, T.~T. (1995).
\newblock Aperiodic stochastic resonance in excitable systems.
\newblock \emph{Physical Review E}, 52(4):R3321.

\bibitem[Faisal et~al.(2008)]{faisal2008noise}
Faisal, A.~A., Selen, L.~P., and Wolpert, D.~M. (2008).
\newblock Noise in the nervous system.
\newblock \emph{Nature Reviews Neuroscience}, 9(4):292--303.

\bibitem[Gammaitoni et~al.(1998)]{gammaitoni1998stochastic}
Gammaitoni, L., H{\"a}nggi, P., Jung, P., and Marchesoni, F. (1998).
\newblock Stochastic resonance.
\newblock \emph{Reviews of Modern Physics}, 70(1):223--287.

\bibitem[Kanerva(2009)]{kanerva2009hyperdimensional}
Kanerva, P. (2009).
\newblock Hyperdimensional computing: An introduction to computing in distributed representation with high-dimensional random vectors.
\newblock \emph{Cognitive Computation}, 1(2):139--159.

\bibitem[Kirkpatrick et~al.(1983)]{kirkpatrick1983optimization}
Kirkpatrick, S., Gelatt, C.~D., and Vecchi, M.~P. (1983).
\newblock Optimization by simulated annealing.
\newblock \emph{Science}, 220(4598):671--680.

\bibitem[Lee et~al.(2013)]{lee2013network}
Lee, U., Ku, S., Noh, G., et~al. (2013).
\newblock Disruption of frontal--parietal communication by ketamine, propofol, and sevoflurane.
\newblock \emph{Anesthesiology}, 118(6):1264--1275.

\bibitem[Mainen and Sejnowski(1995)]{mainen1995reliability}
Mainen, Z.~F. and Sejnowski, T.~J. (1995).
\newblock Reliability of spike timing in neocortical neurons.
\newblock \emph{Science}, 268(5216):1503--1506.

\bibitem[Mashour(2014)]{mashour2014bottom}
Mashour, G.~A. (2014).
\newblock Top-down mechanisms of anesthetic-induced unconsciousness.
\newblock \emph{Frontiers in Systems Neuroscience}, 8:115.

\bibitem[McDonnell and Ward(2011)]{mcdonnell2011benefits}
McDonnell, M.~D. and Ward, L.~M. (2011).
\newblock The benefits of noise in neural systems: bridging theory and experiment.
\newblock \emph{Nature Reviews Neuroscience}, 12(7):415--425.

\bibitem[Moss et~al.(2004)]{moss2004stochastic}
Moss, F., Ward, L.~M., and Sannita, W.~G. (2004).
\newblock Stochastic resonance and sensory information processing: a tutorial and review of application.
\newblock \emph{Clinical Neurophysiology}, 115(2):267--281.

\bibitem[Oizumi et~al.(2014)]{oizumi2014phenomenology}
Oizumi, M., Albantakis, L., and Tononi, G. (2014).
\newblock From the phenomenology to the mechanisms of consciousness: Integrated Information Theory 3.0.
\newblock \emph{PLoS Computational Biology}, 10(5):e1003588.

\bibitem[Plate(1995)]{plate1995holographic}
Plate, T.~A. (1995).
\newblock Holographic reduced representations.
\newblock \emph{IEEE Transactions on Neural Networks}, 6(3):623--641.

\bibitem[Schwitzgebel(2015)]{schwitzgebel2015if}
Schwitzgebel, E. (2015).
\newblock If materialism is true, the United States is probably conscious.
\newblock \emph{Philosophical Studies}, 172(7):1697--1721.

\bibitem[Stocks(2000)]{stocks2000suprathreshold}
Stocks, N.~G. (2000).
\newblock Suprathreshold stochastic resonance in multilevel threshold systems.
\newblock \emph{Physical Review Letters}, 84(11):2310.

\bibitem[Tegmark(2016)]{tegmark2016improved}
Tegmark, M. (2016).
\newblock Improved measures of integrated information.
\newblock \emph{PLoS Computational Biology}, 12(11):e1005123.

\bibitem[Tononi(2004)]{tononi2004information}
Tononi, G. (2004).
\newblock An information integration theory of consciousness.
\newblock \emph{BMC Neuroscience}, 5(1):42.

\bibitem[Tononi(2008)]{tononi2008consciousness}
Tononi, G. (2008).
\newblock Consciousness as integrated information: a provisional manifesto.
\newblock \emph{The Biological Bulletin}, 215(3):216--242.

\bibitem[Tononi et~al.(2015)]{tononi2015integrated}
Tononi, G., Boly, M., Massimini, M., and Koch, C. (2015).
\newblock Integrated information theory: from consciousness to its physical substrate.
\newblock \emph{Nature Reviews Neuroscience}, 17(7):450--461.

\end{thebibliography}

\end{document}
